% Introduction (recto)
\cleardoublepage
\thispagestyle{empty}
\begin{center}
  \textcolor{bookheadingcolor}{\Huge Introduction}\par
\end{center}
\vspace{2em}

\noindent{\large\textbf{The Old Testament Text}}
\vspace{0.5em}

\noindent The Old Testament in this edition is the Septuagint as published by Sir Lancelot Charles Lee
Brenton, originally issued in 1851 by Samuel Bagster \& Sons, London. The digital text was
obtained from ebible.org and was originally prepared for \LaTeX\ by the Brenton LXX \LaTeX\
Print Project (mrgreekgeek).

During the preparation of this edition, extensive OCR and transcription errors were discovered in
the digitized text. To address these, every word of the Septuagint was compared against the
critical editions of Rahlfs (\textit{Septuaginta}, 1935) and Swete (\textit{The Old Testament
in Greek}, 1909--1922) using six automated detection methods. These methods identified vocabulary
errors (words absent from all known editions), character substitution errors, accent and
breathing mark discrepancies, and cases of incorrectly merged or split words.

Through this process, 1,792 corrections were applied across all 53 books of the Old Testament
and the Apocrypha. These corrections fall into two categories:

\textit{OCR corrections} (approximately 1,370): These restore the text to what Brenton's
printed edition actually says. They correct errors introduced during digitization and OCR
processing. Each was verified visually against a PDF of Brenton's \textasciitilde 1900 edition
(\textit{The Septuagint Version of the Old Testament and Apocrypha}, Samuel Bagster, London;
from archive.org).

\textit{Editorial corrections} (approximately 170): These are corrections \textit{away} from
Brenton's text to align with the critical editions---primarily Rahlfs (approximately 146
corrections), with some to Swete (approximately 12), Codex Vaticanus (4), and Codex Alexandrinus
(1). In a small number of cases, an editorial judgment was made based on the available evidence.

Additionally, approximately 2,400 legitimate textual variants were catalogued as valid
differences between Brenton and the critical editions. The complete correction log with full
details is available in the source repository.

\vspace{1em}
\noindent{\large\textbf{The New Testament Text}}
\vspace{0.5em}

\noindent The New Testament text is the Open Greek New Testament (OpenGNT) by Eliran Wong, released under
the Creative Commons Attribution 4.0 International License (CC BY 4.0). The OpenGNT provides an
NA-equivalent text compiled from open resources. No corrections were made to this text.

\vspace{1em}
\noindent{\large\textbf{Versification and Canon}}
\vspace{0.5em}

\noindent The Psalms follow LXX numbering, with Masoretic cross-references provided. The Old Testament
follows the standard Septuagint canonical order, including the deuterocanonical books. The New
Testament follows the standard 27-book ordering.

\vspace{1em}
\noindent{\large\textbf{Source Code}}
\vspace{0.5em}

\noindent The source code used to build this edition is freely available at:\\
github.com/jjorloff1/lxx-nt-greek-bible-builder
